\chapter{Mécanique et théorie des enchères}\label{ch:mx:auction-mechanics-and-theory}

Dans les minutes qui précèdent une enchère de clôture, la bourse publie un prix indicatif et un déséquilibre toutes les quelques secondes, et une grande partie du volume
de l'enchère arrive dans ses dernières secondes. Chaque participant observe les mêmes chiffres et décide du moment où il abat son jeu. L'enchère de clôture fixe les
cours que suivent les fonds indiciels et à partir desquels les fonds sont valorisés~; sa conception, ce qui est publié, le moment où les ordres sont acceptés et l'instant exact où elle prend fin détermine qui
peut peser sur le cours de clôture et qui en supporte le coût. Ce chapitre fait tourner des enchères de clôture dans le simulateur de bourse, mesure la convergence du prix indicatif vers le
prix final, calcule le prix d'un ordre tardif sous une fin fixe puis sous une \omterm{def:mx:auction-mechanics-and-theory:randomend}{fin aléatoire}, et rejoue la course au sniping que les \omterm{def:mx:auction-mechanics-and-theory:fba}{enchères par lots fréquentes} ont été proposées pour arrêter.

\section{Règles du croisement revisitées}

Une enchère (livre~1, chapitre~13) rassemble des ordres sans les apparier et les règle tous à un même prix~: celui qui maximise le
volume exécuté, puis minimise le surplus, puis suit le côté porteur du surplus, puis reste le plus proche d'un prix de référence (le
\texttt{firm.auction} du livre~1). Avec des achats de 300 actions au marché, 200 à 10.02 et 400 à 10.00, et des ventes de 500 à 10.00, 300 à 10.01 et 200 à 10.03, le volume
est de 500 à 10.00, 10.01 et 10.02~; le surplus est de 400 à l'achat à 10.00 et de 300 à la vente sur les deux autres niveaux, et la pression vendeuse retient le plus bas, 10.01.
Madhavan (1992) a comparé les mécanismes~: une enchère périodique offre une meilleure efficience des prix et peut fonctionner là où les mécanismes continus échouent, mais les traders
renoncent à la continuité et supportent des coûts d'information plus élevés.

Le moteur de \texttt{firm.exchsim} fait tourner les enchères~: les ordres sont acceptés sans appariement, les \omterm{def:mx:the-limit-order-book:orders}{ordres au marché} à la clôture et les \omterm{def:mx:the-limit-order-book:orders}{ordres à cours limité} à la clôture (tif C) sont mis de côté, le
moteur publie un cours indicatif et un déséquilibre selon un calendrier (message I du flux), puis déclenche le croisement à l'aide de \texttt{firm.auction}. \texttt{firm.auctionsim}
énonce une conception sous forme de politique et la transforme en configuration du simulateur.

\omcode{code/firm/auctionsim/firm_auctionsim.py}{45}{57}{Une enchère de clôture comme politique~: les phases de l'appel, et un collar comme bande de prix dès le début de l'appel.}\label{lst:mx:auction-mechanics-and-theory:policy}

\section{Publication des cours indicatifs et du déséquilibre}

Le cours indicatif est le prix auquel l'enchère se réglerait si elle se terminait maintenant, et le déséquilibre, les actions qui resteraient d'un seul côté. Publier ces données
invite la contrepartie~: un fort déséquilibre acheteur à un cours supérieur à la valeur est une offre à quiconque va vendre. L'enchère de clôture de ce chapitre dure cinq minutes, avec
un cours indicatif toutes les secondes. Des \omterm{def:mx:the-limit-order-book:orders}{ordres à cours limité} à la clôture permanents, à hauteur de 500 actions par tick, se trouvent de chaque côté de 100.00~; des \omterm{def:mx:the-limit-order-book:orders}{ordres au marché} à la clôture aléatoires, de 100 à 1\,000
actions arrivent cinq fois par minute jusqu'à cinq secondes avant la fin~; un fonds envoie un \omterm{def:mx:the-limit-order-book:orders}{ordre au marché} à la clôture d'achat de 10\,000 actions par minute pendant l'enchère.
Les répondeurs lisent chaque cours indicatif et, une fois sur deux, lorsque le cours a été poussé à plus d'un tick de 100.00, proposent le côté opposé à 100.01 (ou 99.99)
proportionnellement à la distance, à raison de 250 actions par tick jusqu'à 3\,000~: ils fournissent la contrepartie du déséquilibre au cours qu'ils estiment, et non au cours auquel il a été poussé.

\omcode{code/microstructure/10-auction-mechanics-and-theory/python/mx_auction.py}{57}{68}{Un répondeur au déséquilibre publié~: il propose le côté opposé à sa propre valeur augmentée d'un tick, proportionnellement à la poussée.}\label{lst:mx:auction-mechanics-and-theory:resp}

\begin{omfigure}
\begin{tikzpicture}
\begin{groupplot}[group style={group size=2 by 1,horizontal sep=1.6cm},width=6.6cm,height=5cm,
  tick label style={font=\footnotesize},label style={font=\small},grid=major,grid style={gray!25},table/col sep=comma,
  x dir=reverse,xlabel={secondes avant le croisement}]
\nextgroupplot[title={\small prix indicatif, une séance},xmin=0,xmax=300,ymin=99.97,ymax=100.22,
  yticklabel style={/pgf/number format/.cd,fixed,precision=2}]
\addplot[omDef,very thick,const plot] table[x=s,y=price]{figdata/microstructure/10-auction-mechanics-and-theory/path.csv};
\nextgroupplot[title={\small $|$indicatif $-$ final$|$ (pas de cotation)},xmode=log,xmin=0.7,xmax=150,ymin=0,ymax=6]
\addplot[omThm,very thick,mark=*,error bars/.cd,y dir=both,y explicit] table[x=s,y=gap,y error=se]{figdata/microstructure/10-auction-mechanics-and-theory/gap.csv};
\end{groupplot}
\end{tikzpicture}

\omcaption{L'enchère de clôture. À gauche~: le cours indicatif d'une séance simulée~; l'achat du fonds le fait monter, les répondeurs le font redescendre, et les ordres tardifs
fixent le cours de clôture. À droite~: l'écart entre le dernier cours indicatif et le prix final, en fonction du temps restant, moyenne sur 24 séances. Données~:
\texttt{mx\_auction.auction\_study}.}\label{fig:mx:auction-mechanics-and-theory:path}
\end{omfigure}

Sur 24 séances simulées (\cref{fig:mx:auction-mechanics-and-theory:path}), le dernier cours indicatif est à 4.7 \omterm{def:mx:the-limit-order-book:tick}{pas de cotation} du prix final 120 secondes avant le
croisement, 3.5 pas à 60 secondes, 2.2 à 30 et 1.0 à 10~; durant les cinq dernières secondes, lorsque plus aucun ordre n'arrive, il coïncide avec le prix final.
Le volume apparié arrive plus tôt que l'information sur les prix~: 76~\% des actions finalement appariées le sont déjà 120 secondes avant le croisement, 89~\% à 60 secondes et 95~\% à 30 secondes. Le volume est
fixé tôt et le prix tardivement, car celui-ci est fixé par les ordres marginaux.

\begin{dated}{2026-09}{Règles de l'enchère de clôture sur deux bourses}\label{dat:mx:auction-mechanics-and-theory:rules}
Nasdaq publie les informations sur le déséquilibre net du croisement de clôture entre 15 h 50 et 16 h 00 ET~; ses ordres de clôture destinés uniquement à compenser les
déséquilibres, doivent être fixés et ne s'exécuter qu'à un prix égal ou supérieur (achat) ou égal ou inférieur (vente) au meilleur cours acheteur ou au meilleur cours vendeur de 16 h 00 ET. À la London Stock Exchange, chaque appel d'enchère et
chaque prolongation est suivie d'une période aléatoire pouvant aller jusqu'à 30 secondes (le guide de la bourse de mars 2020).
\end{dated}

\section{Heures limites, collars et fins aléatoires}

\begin{definition}[Heure limite de l'enchère]\label{def:mx:auction-mechanics-and-theory:cutoff}
Une \emph{heure limite de l'enchère}\index{heure limite de l'enchère} (auction cut-off time) est le moment au-delà duquel certains ordres ne peuvent plus être saisis, modifiés ni annulés pour
l'enchère, typiquement les \omterm{def:mx:the-limit-order-book:orders}{ordres au marché} à la clôture et les \omterm{def:mx:the-limit-order-book:orders}{ordres à cours limité} à la clôture, tandis que les ordres qui compensent le déséquilibre publié peuvent encore être acceptés.
\end{definition}

\begin{definition}[Collar d'enchère]\label{def:mx:auction-mechanics-and-theory:collar}
Un \emph{collar d'enchère}\index{collar d'enchère} (auction collar) est une fourchette de prix autour d'un prix de référence, hors de laquelle l'enchère ne peut pas se conclure, ou des ordres ne peuvent pas être
fixés~; si le croisement devait tomber hors de cette fourchette, la plateforme prolonge l'appel ou règle au niveau du collar.
\end{definition}

\begin{definition}[Fin aléatoire]\label{def:mx:auction-mechanics-and-theory:randomend}
Une \emph{fin aléatoire}\index{fin aléatoire} (random end) met fin à un appel d'enchère à un instant aléatoire dans une fenêtre connue après l'heure de fin prévue, de sorte que personne ne puisse être certain
d'être le dernier à saisir un ordre.
\end{definition}

Les trois répondent au même problème~: le dernier ordre entré dans un appel peut en déplacer le cours, et personne ne peut y répondre. Une heure limite arrête les ordres qui créeraient
des déséquilibres tardifs~; un collar borne la distance dont un ordre peut déplacer le cours~; une \omterm{def:mx:auction-mechanics-and-theory:randomend}{fin aléatoire} ôte la certitude d'être le dernier. \texttt{firm.auctionsim}
implémente le collar comme une bande de prix dès le début de l'appel (les ordres situés hors de celle-ci sont rejetés, \texttt{Out\_J} avec le motif B), l'heure limite avec le
contrôle de gel du moteur, et la \omterm{def:mx:auction-mechanics-and-theory:randomend}{fin aléatoire} avec l'heure de fin uniforme du simulateur.

\section{Comportement stratégique à proximité du croisement}

Que vaut la dernière seconde~? L'expérience ajoute un \omterm{def:mx:the-limit-order-book:orders}{ordre au marché} à la clôture d'achat supplémentaire de 3\,000 actions à chaque séance, envoyé 60, 10 ou 1 seconde avant la
fin prévue, et mesure de combien il déplace le cours de clôture par rapport à la même séance sans cet ordre (\cref{fig:mx:auction-mechanics-and-theory:late}).
Avec une fin fixe, l'ordre déplace la clôture de 3.0 \omterm{def:mx:the-limit-order-book:tick}{pas de cotation} (erreur type 0.6) s'il est envoyé une minute à l'avance, de 3.5 pas dix secondes à l'avance et de 4.7 pas (0.4) durant la dernière seconde~:
envoyé tard, il consomme l'intérêt affiché sans que les répondeurs aient le temps de proposer. Avec une \omterm{def:mx:auction-mechanics-and-theory:randomend}{fin aléatoire} pouvant aller jusqu'à 30 secondes, l'ordre tardif déplace la clôture de
3.7 \omterm{def:mx:the-limit-order-book:tick}{pas de cotation} (0.6)~: les répondeurs disposent du reste de la fenêtre aléatoire. Le gain tiré d'un envoi en dernier, 1.75 \omterm{def:mx:the-limit-order-book:tick}{pas de cotation} avec une fin fixe, tombe à 0.75 avec une
\omterm{def:mx:auction-mechanics-and-theory:randomend}{fin aléatoire}. Ce gain est ce qu'un participant voulant déplacer la clôture (pour valoriser un portefeuille, ou pour déclencher quelque chose indexé dessus) achète en attendant~; il est aussi
le prix que paie un ordre tardif légitime, ce qui explique que les fournisseurs de liquidité attendent eux aussi.

\begin{omfigure}
\begin{tikzpicture}
\begin{axis}[width=9cm,height=5.2cm,xmode=log,x dir=reverse,xtick={1,10,60},xticklabels={1,10,60},xlabel={\small secondes avant la fin prévue},
  ylabel={\small déplacement de la clôture (pas de cotation)},tick label style={font=\footnotesize},ymin=0,ymax=6,xmin=0.7,xmax=90,grid=major,grid style={gray!25},
  legend columns=2,legend style={font=\footnotesize,at={(0.5,-0.3)},anchor=north,draw=none,fill=none,
  /tikz/every even column/.append style={column sep=8pt}},table/col sep=comma]
\addplot[omDef,very thick,mark=*,error bars/.cd,y dir=both,y explicit] table[x=tau,y=fixed,y error=fixed_se]{figdata/microstructure/10-auction-mechanics-and-theory/late.csv};
\addplot[omThm,very thick,mark=square*,error bars/.cd,y dir=both,y explicit] table[x=tau,y=random,y error=random_se]{figdata/microstructure/10-auction-mechanics-and-theory/late.csv};
\legend{fin fixe,fin aléatoire (jusqu'à 30 s)}
\end{axis}
\end{tikzpicture}

\omcaption{Jusqu'où un achat au marché à la clôture de 3\,000 actions déplace le cours de clôture selon le moment de son envoi, moyenne et erreur type sur 24
séances simulées. La dernière seconde vaut le plus avec une fin fixe. Données~: \texttt{mx\_auction.auction\_study}.}\label{fig:mx:auction-mechanics-and-theory:late}
\end{omfigure}

Les clôtures réelles ajoutent ce que la simulation laisse de côté~: des participants qui choisissent leur propre moment d'intervention, des ordres portant uniquement sur le déséquilibre, et de gros flux d'indices connus à l'avance.
Bogousslavsky et Muravyev (2023) étudient qui négocie à la clôture et ce que cela implique pour la \omterm{def:mx:fragmentation-and-routing:discovery}{formation des prix} et la liquidité.

\section{Enchères par lots fréquentes}

\begin{definition}[Enchère par lots fréquente]\label{def:mx:auction-mechanics-and-theory:fba}
Une \emph{enchère par lots fréquente}\index{enchère par lots fréquente} remplace la confrontation continue des ordres par des enchères à prix uniforme tenues à intervalles fixes et courts~:
les ordres arrivant pendant un intervalle sont traités comme simultanés et exécutés ensemble à la fin de celui-ci.
\end{definition}

Budish, Cramton et Shim (2015) ont soutenu que la course à la vitesse du trading à haute fréquence est un symptôme de la confrontation continue~: lorsqu'une information publique fait bouger le
cours, un trader rapide peut saisir une cotation avant que son propriétaire ne l'annule, et c'est le plus rapide qui gagne. Ils ont proposé des enchères par lots, par exemple toutes les dix
centièmes de seconde, de sorte que le temps devienne discret et qu'un petit avantage de vitesse cesse de décider de qui négocie. Le simulateur les exécute avec
\texttt{batch\_interval\_ns} (\texttt{firm.auctionsim.batch\_venue}).

La course est facile à mettre en scène. Un teneur de cotation affiche un lot d'un tick de part et d'autre d'une référence publique~; la référence saute de deux ticks 400 fois~; à chaque saut, le
teneur annule puis recote après sa latence, et un sniper envoie un ordre pour prendre la cotation périmée après la sienne. Le teneur étant 150 microsecondes plus lent que le sniper, celui-ci prend la cotation
périmée lors de 99.75~\% des sauts sur le marché continu, 11.75~\% avec des lots toutes les millisecondes, 1.25~\% toutes les 10 millisecondes et 0.25~\% toutes les 100. Avec un écart d'une milliseconde, on obtient 99.75~\%, 99.75~\%, 11.0~\% et 0.75~\%. La part est proche du rapport entre l'écart de latence et l'intervalle de lot, plafonné à un~: le sniper ne gagne que lorsqu'un lot se clôt entre son ordre et l'annulation du teneur. Un intervalle de lot plus long que les écarts de latence qui comptent transforme la course en enchère.
10 millisecondes et 0.25~\% toutes les 100. Avec un écart d'une milliseconde, 99.75~\%, 99.75~\%, 11.0~\% et 0.75~\%. La part est proche de l'écart de latence divisé par
l'intervalle entre lots, plafonnée à un~: le sniper ne gagne que si un lot se clôt entre son ordre et l'annulation du coteur. Un intervalle entre lots plus long que les
écarts de latence qui comptent transforme la course en enchère.

\section{Tutoriel~: les dix dernières secondes}\label{tut:mx:auction-mechanics-and-theory:tut}

\begin{tutorial}
\textbf{Objectif.} Lancer des appels de clôture, observer le cours indicatif converger, fixer le prix de la dernière seconde selon deux règles de fin, et rejouer la course au sniping.
\textbf{État final~:} \cref{fig:mx:auction-mechanics-and-theory:path,fig:mx:auction-mechanics-and-theory:late} et les chiffres des sections 2, 4 et 5.
\begin{tutsteps}
\item \textbf{Conception.} \texttt{firm\_auctionsim.ClosingAuction(call\_s, indicative\_s, cutoff\_s, collar, random\_end\_s)}, puis
\texttt{.phases(\dots)} et \texttt{.controls(\dots)} dans \texttt{ExchangeConfig} et \texttt{Simulator.add\_events}.
\item \textbf{Séances.} \texttt{mx\_auction.close\_session(seed, random\_end\_s, extra)}~; \texttt{indicatives(res)} et \texttt{final\_cross(res)}.
\item \textbf{Étude.} \texttt{auction\_study()}~: les trajectoires de l'écart et des arrivées, et le déplacement de la clôture par un ordre tardif, sur 24 graines.
\item \textbf{Course.} \texttt{snipe\_race(interval\_ms, quoter\_us=\dots, sniper\_us=\dots)} et \texttt{race\_study()}~; tracer avec \texttt{fig\_auction.py}.
\end{tutsteps}
\textbf{Pour aller plus loin.} Ajouter une heure limite dix secondes avant la fin, pour les seuls \omterm{def:mx:the-limit-order-book:orders}{ordres au marché} à la clôture, et laisser les répondeurs continuer à proposer~; donner aux répondeurs
un retard de cinq secondes et observer la protection de la \omterm{def:mx:auction-mechanics-and-theory:randomend}{fin aléatoire} s'amenuiser.
\end{tutorial}

\section{Construction~: politiques d'enchère}\label{bld:mx:auction-mechanics-and-theory:auctionsim}

\begin{build}
\bfield{Rôle} Conceptions d'enchère pour le simulateur~: les clôtures des algorithmes de référence des chapitres 16 et 28, les réouvertures après suspension du chapitre~25, et
les stratégies d'enchère du livre~11.
\bfield{Interface} \texttt{ClosingAuction(call\_s, indicative\_s, cutoff\_s, collar, random\_end\_s)} avec \texttt{.phases}, \texttt{.controls}~;
\texttt{batch\_venue(interval\_ms)}~; \texttt{indicatives(res, venue, locate)}, \texttt{final\_cross(res, venue, locate)},
\texttt{gap\_path(ind, final, grid\_s)}.
\bfield{Règles} Le moteur conserve les ordres de l'enchère et procède au croisement par l'intermédiaire de \texttt{firm.auction} du livre~1~; les cours indicatifs ne sont publiés que lorsqu'un volume
se croise~; le collar est une bande de prix à partir du début de la phase d'appel~; la \omterm{def:mx:auction-mechanics-and-theory:randomend}{fin aléatoire} est uniforme sur sa fenêtre.
\bfield{Tests d'acceptation} \texttt{code/firm/auctionsim/tests/}~: un appel scripté dont les cours indicatifs et le croisement sont vérifiés à la main~; une \omterm{def:mx:auction-mechanics-and-theory:randomend}{fin aléatoire} située dans sa
fenêtre~; un ordre hors du collar rejeté~; l'intervalle d'un lot de la plateforme.
\bfield{Pour aller plus loin} Extensions de surveillance des prix lorsque le croisement laisserait le cours hors du collar~; ordres portant uniquement sur le déséquilibre~; un appel de clôture sur plusieurs plateformes avec une seule
prix de référence.
\end{build}

\begin{omsources}
\item A.~Madhavan, ``Trading mechanisms in securities markets'', \emph{Journal of Finance} 47(2), 1992.
\item E.~Budish, P.~Cramton and J.~Shim, ``The high-frequency trading arms race: frequent batch auctions as a market design response'', \emph{Quarterly
Journal of Economics} 130(4), 2015.
\item V.~Bogousslavsky and D.~Muravyev, ``Who trades at the close? Implications for price discovery and liquidity'', \emph{Journal of Financial Markets}
66, 2023.
\item Nasdaq, \emph{The Nasdaq Opening and Closing Crosses}, web page, 2026.
\item London Stock Exchange, \emph{Maintaining orderly markets: circuit breakers explained}, 2020.
\end{omsources}

\section{Exercices}

\begin{exercise}[$\star$]\label{exo:mx:auction-mechanics-and-theory:1}
Achats~: 300 actions au marché, 200 à 10.02, 400 à 10.00. Ventes~: 500 à 10.00, 300 à 10.01, 200 à 10.03. Trouver le prix de croisement, le volume et le
surplus avec les règles du livre~1.
\end{exercise}

\begin{exercise}[$\star$]\label{exo:mx:auction-mechanics-and-theory:2}
Le cours indicatif montre 20\,000 actions appariées à 50.10, avec 3\,000 de plus à l'achat. Qu'est-ce que cela apprend à un fournisseur de liquidité, et à quel prix offrirait-il~?
\end{exercise}

\begin{exercise}[$\star$]\label{exo:mx:auction-mechanics-and-theory:3}
Un sniper est plus rapide qu'un coteur de 150 microsecondes. À quelle fréquence, approximativement, prend-il une cotation périmée avec des enchères par lots toutes les 1 et toutes les 10
millisecondes~? Et avec un écart d'une milliseconde et des lots toutes les 10~?
\end{exercise}

\begin{exercise}[$\star\star$]\label{exo:mx:auction-mechanics-and-theory:4}
L'intérêt en attente est de 500 actions par tick côté vendeur. De combien de ticks un achat au marché à la clôture de 3\,000 actions déplace-t-il la clôture si personne ne répond~? Qu'a
mesuré la simulation pour un ordre de dernière seconde, et pourquoi est-ce moins~?
\end{exercise}

\begin{exercise}[$\star\star$]\label{exo:mx:auction-mechanics-and-theory:5}
Pourquoi le volume apparié est-il connu bien avant le prix dans la phase d'appel simulée (76~\% des actions appariées deux minutes avant le croisement, alors que le cours indicatif est encore
4.7 ticks away)~?
\end{exercise}

\begin{exercise}[$\star\star$]\label{exo:mx:auction-mechanics-and-theory:6}
Expliquer pourquoi une \omterm{def:mx:auction-mechanics-and-theory:randomend}{fin aléatoire} réduit le gain à soumettre en dernier, et ce qu'elle coûte aux autres participants.
\end{exercise}

\begin{exercise}[$\star\star\star$]\label{exo:mx:auction-mechanics-and-theory:7}
\emph{Programmation.} Relancer \texttt{snipe\_race} avec des lots toutes les 1, 10 et 100 millisecondes et un écart de latence de 150 microsecondes, et comparer le taux de sniping
à l'écart divisé par l'intervalle. Où l'approximation échoue-t-elle~?
\end{exercise}

\begin{exercise}[$\star\star\star$]\label{exo:mx:auction-mechanics-and-theory:8}
\emph{Trouver la faille.} «~Le cours de clôture est fixé par tout le marché, donc aucun ordre tardif ne peut à lui seul le déplacer beaucoup.~»
\end{exercise}

\section{Problème~: les dix dernières secondes}

\begin{problem}[{Problème du week-end --- les dix dernières secondes}]\label{pb:mx:auction-mechanics-and-theory:1}
Une bourse se demande s'il faut ajouter une \omterm{def:mx:auction-mechanics-and-theory:randomend}{fin aléatoire} à son enchère de clôture et remplacer la négociation continue par des \omterm{def:mx:auction-mechanics-and-theory:fba}{enchères par lots fréquentes}. Répondre avec les
simulator.

\medskip\noindent\textbf{Partie I --- La phase d'appel.}
\begin{enumerate}
\item Énoncer les règles de croisement et les appliquer à l'exercice~1.
\item Que publie le moteur pendant une phase d'appel, et quand~?
\item Décrire les participants de l'enchère simulée.
\item Comment les répondants fixent-ils le prix de leurs offres, et pourquoi est-ce important~?
\end{enumerate}

\medskip\noindent\textbf{Partie II --- Convergence.}
\begin{enumerate}[resume]
\item Donner l'écart entre le cours indicatif et le cours final à 120, 60, 30 et 10 secondes.
\item Pourquoi est-il nul durant les cinq dernières secondes~?
\item Quelle part du volume apparié est connue deux minutes et une minute avant le croisement~?
\item Que publient ou randomisent le Nasdaq et le London Stock Exchange~?
\end{enumerate}

\medskip\noindent\textbf{Partie III --- La dernière seconde.}
\begin{enumerate}[resume]
\item Définir une limite de cours, un collar et une \omterm{def:mx:auction-mechanics-and-theory:randomend}{fin aléatoire}.
\item De combien un achat de 3\,000 actions déplace-t-il le cours de clôture lorsqu'il est envoyé 60, 10 et 1 seconde avant une fin fixe~?
\item Et avant une \omterm{def:mx:auction-mechanics-and-theory:randomend}{fin aléatoire} pouvant aller jusqu'à 30 secondes~?
\item Quel est le gain d'une soumission en dernier sous chaque règle~?
\item Qui gagne à attendre, et qui paie~?
\end{enumerate}

\medskip\noindent\textbf{Partie IV --- Lots et verdict.}
\begin{enumerate}[resume]
\item Énoncer l'argument de Budish, Cramton et Shim.
\item Décrire la course au sniping et donner ses résultats pour les deux écarts de latence.
\item Dériver l'écart d'approximation divisé par l'intervalle.
\item Que coûte un intervalle d'enchère aux traders qui ne participent pas à la course~?
\item Énoncer le \emph{résultat nommé}~: l'écart entre le cours indicatif et le cours final en fonction du temps restant avant le croisement, et le gain d'une soumission tardive
sous une fin fixe et sous une \omterm{def:mx:auction-mechanics-and-theory:randomend}{fin aléatoire}.
\item La bourse doit-elle ajouter une \omterm{def:mx:auction-mechanics-and-theory:randomend}{fin aléatoire}~? Répondre en une phrase.
\item Doit-elle passer à des lots, et à quel intervalle~?
\end{enumerate}
\end{problem}

\section{Questions d'entretien}

\begin{interviewq}[$\star$ \iqroles{trader}]\label{iq:mx:auction-mechanics-and-theory:1}
Comment le prix d'une enchère de clôture est-il déterminé, et qu'est-ce qu'un déséquilibre~?
\end{interviewq}

\begin{interviewq}[$\star\star$ \iqroles{trader}]\label{iq:mx:auction-mechanics-and-theory:2}
Vous devez acheter 500\,000 actions à la clôture. Quand envoyez-vous l'ordre, et que surveillez-vous~?
\end{interviewq}

\begin{interviewq}[$\star\star$ \iqroles{chercheur}]\label{iq:mx:auction-mechanics-and-theory:3}
Pourquoi certaines bourses terminent-elles les enchères à un moment aléatoire~? Que mesureriez-vous pour vérifier que cela fonctionne~?
\end{interviewq}

\begin{interviewq}[$\star\star$ \iqroles{chercheur}]\label{iq:mx:auction-mechanics-and-theory:4}
Expliquez les \omterm{def:mx:auction-mechanics-and-theory:fba}{enchères par lots fréquentes} et le problème qu'elles résolvent.
\end{interviewq}

\begin{interviewq}[$\star\star$ \iqroles{développeur}]\label{iq:mx:auction-mechanics-and-theory:5}
Implémentez efficacement le cours indicatif d'une enchère, au fil de l'arrivée des ordres. Quelle structure de données conservez-vous~?
\end{interviewq}

\begin{interviewq}[$\star\star\star$ \iqroles{trader, chercheur}]\label{iq:mx:auction-mechanics-and-theory:6}
Un concurrent envoie systématiquement de gros ordres de clôture pendant la dernière seconde. Comment distingueriez-vous une stratégie légitime d'une tentative de déplacer le cours de clôture~?
\end{interviewq}
